% SPDX-License-Identifier: AGPL-3.0-or-later
%
% "Composing the error budget of an optical clock: one rate scalar per
% atom", a short companion paper to the CliffordClock methods-and-validation
% paper (paper/main.tex).
%
% DRAFT ONLY. The project owner is the sole author of record, reviews
% every claim below, and decides whether/when/where to submit this
% document. No agent has submitted, posted, or otherwise distributed
% this paper anywhere.
%
% Every number below that originates from the codebase or from the
% companion validation paper is copied from a committed repository
% file (paper/main.tex, docs/CONVENTIONS.md, or
% notebooks/13_trapped_ion_quantum_motion.ipynb); this paper's own
% build does not recompute any of them. The source is named in the
% prose or a footnote at first use. Build: see paper/composition/Makefile.

\documentclass[aps,pra,superscriptaddress,groupedaddress,nofootinbib]{revtex4-2}

\usepackage{amsmath}
\usepackage{amssymb}
\usepackage{booktabs}
\usepackage{hyperref}
\usepackage{placeins}
\usepackage{orcidlink}

\begin{document}

\title{Composing the error budget of an optical clock: one rate scalar per atom}

\author{Michael Rudolph\,\orcidlink{0000-0002-0495-5723}}
\affiliation{Velar Scientific}
\email{michael.rudolph@velarscientific.com}

\date{\today}

\begin{abstract}
An optical clock's systematic error budget is built today as a table,
one row per physical effect: blackbody radiation, lattice or probe
Stark shifts, the electric-quadrupole shift, motional time dilation,
gravitational redshift. Each row is computed in its own formalism.
The rows are summed linearly, and their uncertainties are combined in
quadrature, the square root of the sum of their squares. Fringe
contrast is measured separately, as an experiment outside the
budget's own rows.
We describe an alternative already implemented and
validated in the open-source CliffordClock pipeline. Every systematic
enters as a multiplicative factor on a single scalar rate, the rate
at which an atom's own proper time advances along its worldline. The
accumulated phase is the time integral of that rate. The clock's
reported fractional shift is the mean of the resulting ensemble phase
distribution. The Ramsey fringe visibility is the modulus of the
mean phasor, fixed by that same distribution's phase variance
under the standard Gaussian-phase closure, the assumption that the
ensemble's accumulated phases are Gaussian-distributed. The additive
table is the first-order term of this product's own expansion. We
give the pipeline's own bounds on the terms it drops: a
DC-Stark$\times$blackbody product correction of order $10^{-33}$, and a
gravitational$\times$field-driven cross term of order $10^{-31}$.
Both bounds sit at least twelve orders of magnitude below the $10^{-18}$-$10^{-19}$
systematic floor of the clocks the budget describes. The value of the
product form is structural. It changes no published total. We show three consequences of the
composed form, using numbers already validated in the companion
paper. First, a multi-surface thermal-environment term needs its
higher-order moments evaluated per surface, before they are combined.
Second, a trapped-ion secular-motion reconciliation composes two
systematics as correlated factors on one rate, closing a discrepancy
of $14\sigma$ down to $0.08\sigma$ against a published budget. Third,
a Ramsey visibility comes from the same worldline ensemble that sets
the mean shift, with no separate fit required. Relativistic timing
already composes rate scalars for its own gravitational and kinematic
terms. This paper's claim is one of scope. It extends the same
composition to the atomic-environment terms of a clock's full budget,
evaluated per quantum state. Fringe contrast comes out of the phase
variance of that same object. The implementation is open and checked
against published budgets. Code,
tests, and the six-line composition rule are public under the
AGPLv3.
\end{abstract}

\maketitle

\section{Introduction}
\label{sec:intro}

A modern optical clock's systematic-uncertainty budget is built one row
at a time, and each row draws on its own kind of data. The
blackbody-radiation shift comes from a Planck-integral model of the
clock's radiative environment. The lattice or probe AC-Stark shift
comes from the trap-light polarizability. The DC-Stark shift and its
spatial gradient come from a field survey. The electric-quadrupole
shift, present for an ion with $J\ge1$, comes from the local
field-gradient tensor and an atomic theory coefficient. The motional
time-dilation shift comes from sideband-thermometry-measured mode
populations. The gravitational redshift comes from a height survey
against a local reference~\cite{Ludlow2015}. Each row is a physically
independent calculation, in its own formalism, against its own inputs.
The rows are then summed linearly to a total fractional shift, and their
uncertainties are combined in quadrature to a total systematic
uncertainty. The clock's fringe contrast, the visibility that sets
how well the interrogation resolves the transition, is
measured separately, as an experimental observation outside the
budget's own rows.

Two recent publications make this practice concrete with real numbers.
Aeppli et al.\ report an $^{87}$Sr optical lattice clock's blackbody
shift as $-4.841720\times10^{-15}$, with a stated uncertainty of
$\pm7.3\times10^{-19}$, from a Planck-integral model of the clock's
measured in-vacuum temperature combined with fitted dynamic
polarizability coefficients~\cite{Aeppli2024}. Their total systematic
uncertainty, after every such row is combined, is
$8\times10^{-19}$~\cite{Aeppli2024}. Bothwell et al., on the same
platform, publish a DC-Stark gradient row of
$+0.3(0.2)\times10^{-20}$ per millimetre from a field survey. In
the same table, a gravitational-redshift slope from two independent
analysis methods gives $-9.8(2.3)\times10^{-20}$ and
$-12.8(2.7)\times10^{-20}$ per millimetre respectively. The two
methods differ from each other by $0.85\sigma$, even though both
describe the same physical effect~\cite{Bothwell2022}. On a different
platform and a different species, Brewer et al.\ report an
$^{27}$Al$^+$/$^{25}$Mg$^+$ trapped-ion clock's own motional
time-dilation row as $-17.3(2.9)\times10^{-19}$, one line among the
rows that sum to their clock's total systematic uncertainty, published
as below $10^{-18}$~\cite{Brewer2019}. A later evaluation of the same
species pair on a different trap, Marshall et al., reports the same
row at $-114.6(3.8)\times10^{-19}$ and a total systematic uncertainty
of $5.5\times10^{-19}$~\cite{Marshall2025}. Four papers, two atomic
species, two very different orders of magnitude for nominally the same
physical row: the practice that produces these numbers is precise and
well established. It is a practice, and this paper asks what one
composed object, replacing the table of separately computed rows,
would look like, and where it would matter.

CliffordClock is an open-source pipeline, described and validated
against published clock budgets in a companion
paper,\footnote{The companion paper is \texttt{paper/main.tex} in the
project repository, \url{https://github.com/velar-mbr/CliffordClock}.
Every case-study number in this paper is copied from that paper, from
\texttt{docs/CONVENTIONS.md}, or from
\texttt{notebooks/13\_trapped\_ion\_quantum\_motion.ipynb}; this
paper's own text does not recompute any of them.} that already implements the alternative this paper
describes. Every systematic in that pipeline, DC-Stark, blackbody
radiation, the electric-quadrupole shift, motional time dilation,
gravitational redshift, enters as a multiplicative factor on one
scalar rate at which an atom's own proper time advances. The
clock's reported shift and its fringe visibility both come out of the
same accumulated-phase object. Section
\ref{sec:law} states that composition rule. Section
\ref{sec:equivalence} derives the additive table as its first-order
expansion and gives the pipeline's own bound on what that expansion
drops. Section \ref{sec:payoff} shows three places, already validated
against published numbers, where the fuller composition changes what
a lab would conclude. Section \ref{sec:prior} places this against
relativistic-timing practice, which already composes rate factors for
gravitational and kinematic terms. Section \ref{sec:close} gives the
composition rule in six lines of code.

\section{The composition law}
\label{sec:law}

This section states the composition rule, and shows how the
clock's reported shift and its fringe visibility both come from the
same object. Write $\nu(t)$ for the clock transition's instantaneous frequency
along one atom's worldline, its own path through space, time, and
internal quantum state. In the standard treatment, each systematic
$k$ contributes an independently computed fractional perturbation
$p_k$, and the rows are added: $\nu(t)/\nu_0 - 1 \approx \sum_k p_k$.
The composition rule this paper describes instead multiplies:
\begin{equation}
  \nu(t) \;=\; \nu_0 \prod_k \bigl(1 + p_k(x(t))\bigr) ,
  \label{eq:rate}
\end{equation}
where $x(t)$ is the atom's full state along its worldline at time
$t$, position, velocity, and any internal or motional quantum number
each term's own physics depends on. Equation~\eqref{eq:rate} is the
local proper-time rate: general relativity already has names for a
scalar factor multiplying an ideal clock's rate this way, the lapse
function in a $3{+}1$ decomposition of spacetime, or, for a static
observer, the redshift factor $\sqrt{-g_{00}}$. Each term $p_k$ is
still computed by its own physics. The DC-Stark shift comes from the
transition's differential polarizability and the local field. The
blackbody shift comes from a Planck-integral model of the radiative
environment. The quadrupole shift comes from the local field-gradient
tensor. The motional shift comes from the atom's own velocity or
vibrational occupation. The gravitational term comes from height in a
local potential.
Nothing about any individual term changes. What changes is how the
terms are put back together. Equation~\eqref{eq:rate} multiplies them
as factors on one shared rate. A budget table instead adds them, one
row per systematic, an addition Sec.~\ref{sec:equivalence} shows is
this product's own first-order truncation.

An atom's accumulated phase over an interrogation of proper duration
$T$ is the time integral of that rate,
\begin{equation}
  \Delta\Phi \;=\; 2\pi \int_0^T \nu(t)\,dt ,
  \label{eq:phase}
\end{equation}
in radians. An ensemble of $M$ atoms, each with its own worldline
(its own trajectory, motional state, and local field environment),
gives $M$ accumulated phases $\Delta\Phi_1,\dots,\Delta\Phi_M$. The
clock's reported fractional frequency shift is the ensemble mean of
these phases, normalized by the unperturbed phase $2\pi\nu_0 T$:
\begin{equation}
  \frac{\Delta\nu}{\nu_0} \;=\; \frac{\langle\Delta\Phi\rangle}{2\pi\nu_0 T} - 1 ,
  \label{eq:shift}
\end{equation}
the first moment of the ensemble's phase distribution, converted to a
fractional frequency. The Ramsey fringe visibility, the same
contrast a lab measures at the end of an interrogation, is the
modulus of the mean phasor over that same ensemble,
\begin{equation}
  V \;=\; \bigl|\langle e^{i\Delta\Phi}\rangle\bigr| \;\le\; 1 .
  \label{eq:visibility}
\end{equation}
For phases that are Gaussian-distributed across the ensemble, the
characteristic-function identity $\langle e^{i\Delta\Phi}\rangle =
e^{i\langle\Delta\Phi\rangle - \sigma_\Phi^2/2}$ gives the closed form
$V = \exp(-\sigma_\Phi^2/2)$ exactly. The visibility is fixed by the
phase distribution's variance, the same distribution whose mean
fixed the shift in Eq.~\eqref{eq:shift}. One worldline ensemble, run once, yields both
numbers. Section~\ref{sec:visibility} validates this identity against
a live pipeline run.

Two constructions of the ensemble average must be told apart, because
one of them silently destroys the visibility signal. Averaging the
per-worldline phases first, and only then taking the phasor,
$|\exp(i\langle\Delta\Phi\rangle)|$, gives modulus $1$ for any spread
in the underlying phases. That happens because the spread is
discarded before the modulus is ever taken. The correct order,
Eq.~\eqref{eq:visibility},
sums the phasors themselves and takes the modulus of that sum, which
is what lets a real phase spread show up as $V<1$. Both orderings are
encoded as opposing kill tests in the pipeline's own test suite
(\texttt{tests/test\_coherent\_visibility.py}), so a future change that
collapses the two back together is caught mechanically.

\section{The additive table as the first-order expansion}
\label{sec:equivalence}

The standard additive budget turns out to be the first term of the
product rule's own expansion, and every higher-order term it drops is
small enough to ignore at current clock precision. Expanding the
product in Eq.~\eqref{eq:rate} to first order in each
$p_k$ recovers the additive budget exactly:
\begin{equation}
  \prod_k (1+p_k) \;=\; 1 + \sum_k p_k \;+\; \sum_{j<k} p_j p_k \;+\; \dots
  \label{eq:expansion}
\end{equation}
The standard table, one row per systematic, rows summed linearly, is
the truncation of Eq.~\eqref{eq:expansion} at first order. Every
cross term the truncation drops is a product of two or more
systematics that are each already small, and at the field's current
performance those products are far smaller still. The pipeline's own
composition rules state each bound directly, as part of the
codebase's own documentation. The
DC-Stark and blackbody-radiation terms are independent scalar
perturbations of the atom's internal energy. Their second-order cross
term averages to exactly zero for an isotropic thermal field, an
exact cancellation. Their multiplicative product correction, the
Stark$\times$BBR cross term of Eq.~\eqref{eq:expansion}, is bounded
at order $10^{-33}$.\footnote{docs/CONVENTIONS.md, section 13,
``(E33) Scalar pivot composition.''} The gravitational-redshift term
depends only on an atom's position in the local potential, and never on
its motional or internal state, so its cross term with any
field-driven shift is bounded at order $10^{-31}$.\footnote{Companion
paper, Sec.~II.E, ``The gravitational-redshift pivot term and extended
samples.''} Both bounds sit at least twelve orders of magnitude below the
$10^{-18}$-to-$10^{-19}$ systematic floor the clocks of
Sec.~\ref{sec:intro} report.

This comparison settles what the product form is for. It states one
composition rule that explains why the additive table works as well
as it does. It also gives that rule's own quantitative boundary, the
order at which it would stop working, in the same units the rows
themselves are reported in. The additive budgets of Sec.~\ref{sec:intro}
are already correct to a precision no experiment now approaches. This
paper does not correct any of those published totals, and revises no
published number. The three
cases of Sec.~\ref{sec:payoff} show something different: where the
product form's other structural consequences change a conclusion. A
nonlinear function of a distributed quantity gets evaluated before it
is averaged; two systematics compose as correlated factors on one
shared rate; a fringe visibility gets read off the same object that
sets the
shift. None of the cross-term corrections above ever enters that
arithmetic.

\section{Three consequences of composing on one rate}
\label{sec:payoff}

\subsection{Per-state evaluation before averaging: a multi-surface thermal environment}
\label{sec:bbr}

A blackbody-radiation budget row is ordinarily a single number: a
lab characterizes its enclosure by one effective temperature and
evaluates a Planck-integral shift formula at that one temperature. A
clock deployed outside a temperature-stabilized laboratory chamber,
in a vehicle or behind a viewport a diurnal cycle warms, sees an
enclosure whose visible surfaces can differ by tens of kelvin. The
blackbody shift's dynamic term scales as $T^6$ through $T^{10}$,
higher powers than the $T^4$ power the static shift and the
naturally chosen effective temperature both track. A single
temperature can be chosen to match the static ($T^4$) moment of a
non-uniform environment. That choice generally misses the
environment's higher dynamic moments, and the companion paper
quantifies the resulting mismatch. It crosses $10^{-18}$,
comparable to a modern lattice clock's total systematic budget, at
roughly an $11\,$K spread across two comparably sized surfaces, and
at roughly $19\,$K for a small warm viewport against a larger chamber
body.\footnote{Companion paper, Sec.~II.D, ``The multi-surface
thermal environment.''}

The pipeline's multi-surface extension evaluates
the blackbody shift's moments per surface, before they are combined
into a single number.
Writing $w_i$ for each surface's solid-angle fraction and $T_i$ for
its temperature, the environment's $n$-th weighted moment is
$M_n = \sum_i w_i (T_i/T_0)^n$. The static and dynamic terms of
the shift formula are each evaluated against their own moment $M_n$,
in place of a single power $(T/T_0)^n$ of one effective
temperature.\footnote{docs/CONVENTIONS.md, section 13, ``(E37)
Multi-surface thermal environment.''} A uniform, single-surface
environment reduces this construction to the ordinary single-$T$
formula term for term: the two paths share the same
coefficient-evaluation code and agree at floating-point precision.
The construction is validated
against a real transportable clock's own position-resolved
measurement. Replaying Nosske et al.'s published aperture geometry,
interior emissivity, and shield temperatures through the pipeline's
enclosure-and-aperture model predicts a differential shift, as the
atoms move from the shield's center out through an aperture, of
$-3.325\times10^{-15}$. Nosske et al.'s own closed-form model gives
$-3.32(7)\times10^{-15}$ for that same shift, and their own
measurement gives $-3.33(3)\times10^{-15}$. The pipeline's prediction
lands within $5.011\times10^{-18}$ of
their model and $4.989\times10^{-18}$ of their measurement, well
inside either one's own stated uncertainty.\footnote{Companion paper,
Sec.~IV.J, ``The PTB position-resolved thermal-environment
validation,'' citing Nosske et al.~\cite{Nosske2025}.}

The point generalizes past blackbody radiation. A single effective
value often stands in for a physically distributed quantity
throughout an optical-clock budget: an effective lattice depth over a
trapped ensemble's spatial extent, an effective field over a cloud's
own spread. Each substitution risks the same failure mode a single
effective temperature shows here. A nonlinear function of the
distributed quantity, evaluated at one collapsed value, does not
equal that same nonlinear function's own moments evaluated across the
real distribution and then combined. The composition rule of
Sec.~\ref{sec:law} states the same discipline once, for every
systematic and every ensemble the pipeline evaluates. The
distribution comes first. The moments, the mean that sets the
shift and the variance that sets the visibility, come from
integrating over the full distribution, before any collapse to one
representative value.

\subsection{Two correlated systematics as factors on one rate: reconciling a trapped-ion secular-motion budget}
\label{sec:motional}

A trapped two-ion clock's motional time-dilation shift depends on two
physically distinct pieces of the same underlying motion. One piece
is how much of each shared vibrational mode the clock ion itself
carries, its participation. The other is how much extra oscillation,
at the trap's own radio-frequency drive, that secular motion itself
induces, its intrinsic micromotion. A budget table built one row at a
time can compute either piece on its own. Marshall et al.'s published
$^{27}$Al$^+$/$^{25}$Mg$^+$ two-ion crystal is a real case where doing
so goes badly wrong before it goes right. Table~\ref{tab:motional}
reproduces the companion paper's own six-step progression through this
reconciliation. Each row is built from Marshall et al.'s own
published six-mode frequencies and occupations, and reported against
Marshall et al.'s own published secular-motion band of
$-114.6(3.8)\times10^{-19}$.\footnote{Companion paper,
Secs.~IV.H-I, ``The Al-ion secular-motion arithmetic reproduction''
and ``Per-mode participation and the intrinsic-micromotion
enhancement''; docs/CONVENTIONS.md section 16.}

\begin{table}[t]
\centering
\caption{Reconciling the published Al$^+$ secular-motion total: the
pipeline's full treatment, the bold final step, lands inside Marshall
et al.'s published band at $0.08\sigma$. The five steps above it are
deliberately incomplete treatments of the same published six-mode
inputs~\cite{Marshall2025}, read top to bottom, each adding one piece
of physics the previous step omitted; their verdicts are kept to show
what each omission costs, and step~1's agreement is the accidental
near-cancellation the text explains. Copied from the companion
paper's own progression table and Fig.~8 there. $\sigma$ is the
distance from Marshall et al.'s published band, combining the
published and predicted uncertainties in quadrature. Participation
and micromotion enter as two multiplicative factors on the same
per-mode rate (Eq.~\eqref{eq:rate}).}
\label{tab:motional}
\begin{tabular}{clcc}
\toprule
Step & Treatment & $\sigma$ & Verdict \\
\midrule
1 & Single species mass, every mode & 0.10 & MET \\
2 & Participation, mass-ratio closed form & 14.11 & NOT MET \\
3 & Participation, measured-spectrum fit & 14.01 & NOT MET \\
4 & Participation $\times$ leading-order micromotion & 1.62 & NOT MET \\
5 & Participation $\times$ exact-Floquet micromotion & 1.62 & NOT MET \\
\textbf{6} & \textbf{Coupled two-ion Floquet, constrained fit} & \textbf{0.08} & \textbf{MET} \\
\bottomrule
\end{tabular}
\end{table}

Step one carries one species mass through every one of the six
modes, with no participation split and no micromotion correction. It
lands inside Marshall et al.'s own published band at $0.10\sigma$.
That agreement is a near-cancellation of two large omissions moving
in opposite directions. The true clock-ion participation in a shared
mode is close to one-half. The true intrinsic-micromotion
enhancement on the radial modes is close to a factor of two. The
single-mass step's implicit choice of one for both factors happens to
land close to their true product. That is an accident of the
approximation, one the next two steps expose. Supplying only the
participation correction, steps two and three, removes that
accidental cancellation without replacing it. The resulting total
under-predicts the published band by a factor of close to two, a
$14\sigma$ discrepancy against Marshall et al.'s own stated
uncertainty. Composing the intrinsic-micromotion enhancement back in
as a second multiplicative factor on the same per-mode rate, steps
four and five, closes most of that distance, to $1.62\sigma$. That
result still falls outside the published band. The
remaining gap traces to a specific mechanism the companion paper
identifies: the two-mode pairs sharing an axis run at
Coulomb-shifted collective frequencies a single-ion enhancement
factor cannot see. Solving the two ions' coupled equations of
motion, then fitting the resulting two-parameter trap model
to all four measured radial frequencies at once, closes the budget to
$0.08\sigma$, step six, comfortably inside Marshall et al.'s own
published band.

A budget table computes participation and micromotion as separate
rows, or omits the smaller of the two once a single-mass evaluation
already looks close to the published value. Composing them as two
correlated factors on the same per-mode rate exposes the
near-cancellation that single-mass evaluation hides. It also lets the
coupled two-ion physics that resolves it enter as one more correction
to that shared rate, the same object every other correction in this
section already multiplies.

\subsection{Shift and fringe visibility from the same phase distribution}
\label{sec:visibility}

Section~\ref{sec:law} states that the same worldline ensemble that
sets the mean fractional shift, Eq.~\eqref{eq:shift}, also sets the
Ramsey fringe visibility, Eq.~\eqref{eq:visibility}. The visibility
calculation uses the same accumulated-phase draw as the shift
calculation, with no free parameter of its own. The
companion paper validates the Gaussian closure $V =
\exp(-\sigma_\Phi^2/2)$ this identity predicts against a live
pipeline run. The run sweeps a classical worldline ensemble's
squeezing parameter over several values, at $4000$ worldlines each.
At each value, the pipeline's own reported visibility is compared
against the same run's independently computed Gaussian-closure check.
The largest relative deviation between the two, at the largest
squeezing swept, lands at roughly $0.65\%$. That deviation is
consistent with the $O(\sigma^4/N)$ Monte Carlo scaling the
pipeline's own test-tolerance derivation documents for this
comparison.\footnote{notebooks/13\_trapped\_ion\_quantum\_motion.ipynb,
Sec.~5, ``Ramsey visibility under squeezing''; companion paper,
Sec.~II.H, ``Coherent Ramsey visibility,'' Eq.~(E39).} That
number comes from the same accumulated-phase distribution
Eq.~\eqref{eq:shift} already builds to report the mean shift, with
the same worldline draw and the same weights supplying both outputs.

\section{Relation to prior practice}
\label{sec:prior}

Relativistic timing already composes a rate scalar for a clock's own
systematics, for the terms general relativity itself supplies.
Ashby's treatment of relativity in the Global Positioning System
expresses each GPS satellite clock's rate, relative to a ground
clock, as a product of the gravitational potential term and the
special-relativistic velocity term. The two compose multiplicatively
along the satellite's worldline, the same proper-time-rate object
this paper composes for an optical clock's atomic-environment
terms~\cite{Ashby2003}. Chronometric geodesy extends the same idea in
the other direction. An optical clock's own measured rate, relative
to a reference, is read as a measurement of the
gravitational potential difference between the clock and the
reference. A recent review of atomic clocks for geodesy
surveys this practice across a range of clock platforms and baselines,
each one treating the clock's fractional frequency as a rate factor tied to
the local metric~\cite{Mehlstaubler2018}.

Both of these established practices compose rate scalars for the
gravitational and kinematic terms general relativity itself supplies.
This paper's own claim is about scope. What
CliffordClock adds is the same composition, extended to a clock's
full budget. That budget includes atomic-environment terms with no
counterpart in a GPS satellite's or a geodetic baseline's own rate
factor: the DC-Stark shift, blackbody radiation, and the
electric-quadrupole shift. Those terms are evaluated per quantum state where the
physics calls for it, the per-surface thermal moments of
Sec.~\ref{sec:bbr} and the per-mode motional factors of
Sec.~\ref{sec:motional} among them. Fringe contrast comes out of
that same object, Sec.~\ref{sec:visibility}, as a byproduct of the
shift calculation. The implementation is open, and every
claim above is checked against a published budget in the repository's
own test suite. No claim is made that this is the only way to compose a
clock's systematics, or the first time a rate scalar has been
composed for any of its individual terms. The claim that is made is
about bookkeeping. One rate scalar per atom carries the whole budget:
every effect enters it the same way, and adding an effect means adding
one factor. Effects that share the same underlying motion stay
correlated, because they multiply the same per-mode rate. That
correlation is what closed the budget of Sec.~\ref{sec:motional}. The
shift and the fringe contrast come out of one calculation. The full rule fits in the
listing of the next section, short enough to check by reading it.

\section{The composition, in six lines}
\label{sec:close}

Every equation of Sec.~\ref{sec:law}, the multiplicative rate, the
phase integral, the shift as a mean, and the visibility as a
modulus of the mean phasor, is this short:

\begin{widetext}
\noindent\begin{minipage}{\linewidth}
\footnotesize
\begin{verbatim}
def clock_phase(worldline, terms, nu_0):
    rate = lambda t: nu_0 * prod(1.0 + p(worldline.state(t)) for p in terms)
    return 2.0 * pi * integrate(rate, worldline.proper_time)

phases = [clock_phase(w, terms, NU_0) for w in ensemble]
shift = mean(phases) / (2.0 * pi * NU_0 * T) - 1.0
visibility = abs(mean(exp(1j * array(phases))))
\end{verbatim}
\normalsize
\end{minipage}
\end{widetext}

Each entry of \texttt{terms} is one systematic's own physics, computed
from its own formalism, as any published budget row is; the
composition is everything else in the listing. The source code,
test suite, and the companion validation paper underlying every number
in this paper are available at
\url{https://github.com/velar-mbr/CliffordClock} under the GNU
Affero General Public License v3.0 or later (AGPLv3).

\bibliography{refs}

\end{document}
